\documentclass[a4paper]{article}
% Español (México) con XeLaTeX: fontspec para fuentes Unicode, polyglossia
% para la división silábica y los nombres del documento en la variante mexicana.
\usepackage{fontspec}
\usepackage{polyglossia}
\setdefaultlanguage[variant=mexican]{spanish}
% Los nombres chinos van como «pinyin (original)»: xeCJK da a los caracteres
% Han su propia fuente; sin ella XeLaTeX los omite con un aviso.
% FandolSong no cubre algunos caracteres de nombres (U+73FA); WenQuanYi Zen Hei sí.
\usepackage[AutoFallBack=true]{xeCJK}
\setCJKmainfont{FandolSong-Regular.otf}
\setCJKfallbackfamilyfont{\CJKrmdefault}{WenQuanYi Zen Hei}
% polyglossia no traduce los nombres de listings.
\AtBeginDocument{\providecommand{\lstlistingname}{}\renewcommand{\lstlistingname}{Listado}%
  \providecommand{\lstlistlistingname}{}\renewcommand{\lstlistlistingname}{Índice de listados}}
\usepackage{amsmath,amssymb}
\usepackage{graphicx}
\usepackage[margin=2.35cm]{geometry}
\usepackage[most]{tcolorbox}
\usepackage{etoolbox}
\usepackage{listings}
\usepackage{xcolor}
\usepackage{hyperref}
\usepackage{booktabs,longtable,tabularx,array}
\usepackage{subcaption}
\usepackage{float}
\usepackage{enumitem}
\usepackage{microtype}
\usepackage{fancyhdr}
\usepackage{needspace}

\definecolor{kbBack}{HTML}{F2F6FA}
\definecolor{kbFrame}{HTML}{9DB4CC}
\definecolor{kbTitle}{HTML}{52708F}
\definecolor{ibBack}{HTML}{FBF5E7}
\definecolor{ibFrame}{HTML}{C9A961}
\definecolor{ibTitle}{HTML}{7D5F1E}
\definecolor{wbBack}{HTML}{FCF2F0}
\definecolor{wbFrame}{HTML}{D6A096}
\definecolor{wbTitle}{HTML}{9E4F43}
\definecolor{dbBack}{HTML}{F4F7F3}
\definecolor{dbFrame}{HTML}{AEC2AC}
\definecolor{dbTitle}{HTML}{5F7F64}

\tcbset{softbox/.style={
  enhanced,breakable,boxrule=0.8pt,arc=2pt,left=3mm,right=3mm,
  top=2mm,bottom=2mm,coltitle=white,fonttitle=\bfseries,
  toptitle=0.6mm,bottomtitle=0.6mm
}}
\newtcolorbox{knowledgebox}[1]{softbox,colback=kbBack,colframe=kbFrame,colbacktitle=kbTitle,title={#1}}
\newtcolorbox{importantbox}[1]{softbox,colback=ibBack,colframe=ibFrame,colbacktitle=ibTitle,title={#1}}
\newtcolorbox{warningbox}[1]{softbox,colback=wbBack,colframe=wbFrame,colbacktitle=wbTitle,title={#1}}
\newtcolorbox{dialoguebox}[1]{softbox,colback=dbBack,colframe=dbFrame,colbacktitle=dbTitle,title={#1}}

\lstset{
  language=Python,basicstyle=\ttfamily\small,
  keywordstyle=\color{kbTitle}\bfseries,stringstyle=\color{wbTitle},
  commentstyle=\color{dbTitle},breaklines=true,frame=single,numbers=left,
  numberstyle=\tiny\color{gray},captionpos=b,columns=fullflexible,
  keepspaces=true,showstringspaces=false,extendedchars=true
}
\BeforeBeginEnvironment{lstlisting}{\Needspace{12\baselineskip}}

\hypersetup{colorlinks=true,linkcolor=kbTitle,urlcolor=kbTitle,citecolor=wbTitle}
\setlist{itemsep=0.25em,topsep=0.4em}
\setlength{\parindent}{2em}
\setlength{\parskip}{0.25em}
\newcolumntype{L}[1]{>{\raggedright\arraybackslash}p{#1}}

\providecommand{\notetitle}{Notas del curso: [tema del curso]}
\providecommand{\noteauthors}{Elaboradas a partir de los materiales públicos de Stanford CS25}
\providecommand{\notedate}{Versión reescrita en 2026}
\providecommand{\videochannel}{Stanford Online}
\providecommand{\videopublishdate}{2022}
\providecommand{\videoduration}{[duración del video]}
\providecommand{\videourl}{}
\providecommand{\videocoverpath}{cover.jpg}
\providecommand{\coursesite}{https://web.stanford.edu/class/cs25/}
\providecommand{\repourl}{https://github.com/hqhq1025/ai-course-notes}
\providecommand{\skillversion}{wdkns/wdkns-skills@39f1a04}

\newcommand{\figsource}[1]{\par\vspace{-0.3em}{\footnotesize\color{gray}\emph{Fuente: #1}}\par}
\newcommand{\teachervoice}[1]{\begin{knowledgebox}{Nota de clase}#1\end{knowledgebox}}
\newcommand{\term}[1]{\textbf{#1}}

\pagestyle{fancy}
\fancyhf{}
\fancyfoot[C]{\thepage}
\fancyfoot[R]{\footnotesize\color{gray}\href{\repourl}{AI Course Notes}}
\renewcommand{\headrulewidth}{0pt}
\renewcommand{\footrulewidth}{0pt}

\newcommand{\makecscover}{
\begin{titlepage}
\centering
\vspace*{0.7cm}
{\Large Stanford CS25 · Transformers United\par}
\vspace{0.8cm}
{\huge\bfseries \notetitle\par}
\vspace{0.6cm}
{\large \noteauthors\par}
\vspace{0.25cm}
{\large \notedate\par}
\vspace{0.9cm}
\ifdefempty{\videocoverpath}{}{
  \includegraphics[width=0.86\textwidth,height=0.43\textheight,keepaspectratio]{\videocoverpath}\par
}
\vfill
\begin{tcolorbox}[width=0.92\textwidth,colback=black!2!white,colframe=black!45,boxrule=0.7pt,arc=2pt]
\textbf{Autor o canal del video}: \videochannel\par
\textbf{Fecha de publicación}: \videopublishdate\par
\textbf{Duración del video}: \videoduration\par
\textbf{Sitio del curso}: \href{\coursesite}{\nolinkurl{\coursesite}}\par
\ifdefempty{\videourl}{}{\textbf{Enlace del video}: \href{\videourl}{\nolinkurl{\videourl}}\par}
\textbf{Normas de elaboración}: \skillversion\ + las normas source-first / teacher-voice / visual-QA de este repositorio
\end{tcolorbox}
\end{titlepage}
\tableofcontents
\newpage
}

\usepackage{multicol}

\renewcommand{\notetitle}{Lecture 28: Aligning Open Language Models}
\renewcommand{\noteauthors}{Elaborado a partir de la clase de Nathan Lambert, el deck oficial de 77 páginas y los subtítulos hechos por personas}
\renewcommand{\notedate}{CS25 V4 · Versión reescrita en 2026 con enfoque source-first}
\renewcommand{\videochannel}{Stanford Online}
\renewcommand{\videopublishdate}{Clase: 2024-04-18; publicación: 2024-05-10}
\renewcommand{\videoduration}{1:16:21}
\renewcommand{\videourl}{https://www.youtube.com/watch?v=AdLgPmcrXwQ}
\renewcommand{\videocoverpath}{cover.jpg}

\newcommand{\lecturefigure}[3]{%
\begin{figure}[H]
  \centering
  \includegraphics[width=0.97\textwidth,height=0.49\textheight,keepaspectratio]{slides-images/#1}
  \caption{#2}
  \figsource{#3}
\end{figure}
}

\let\standardtableofcontents\tableofcontents
\renewcommand{\tableofcontents}{%
  \begingroup
  \small
  \setlength{\parskip}{0pt}
  \standardtableofcontents
  \endgroup
}

\begin{document}
\makecscover

\section{Enfoque de la clase: la alineación abierta no es un algoritmo, sino la historia de un sistema}

Esta clase no es un ``RLHF 101''. Nathan Lambert busca responder lo siguiente: después de ChatGPT, ¿cómo pasó la comunidad abierta de base models que solo saben continuar texto a modelos capaces de seguir instrucciones, conversar, negarse a responder y recibir optimización de preferencias? ¿Qué datos, evaluaciones y métodos de entrenamiento cambiaron realmente el ecosistema y qué nombres fueron solo etiquetas de una etapa? Al leer esta clase, conviene ver los modelos, los datos, los algoritmos, el hardware y la evaluación como un sistema de restricciones mutuas, en lugar de buscar una única «receta de alineación».

\subsection{Primero corregir la fuente y luego entender el alcance del curso}

Esta sección aclara primero los límites del material. El ponente oficial de la Lecture 28 es Nathan Lambert, de AI2, y el tema es la alineación de modelos de lenguaje abiertos; el directorio anterior usó por error el deck sobre la historia de las arquitecturas de la clase previa, impartida por Hyung Won Chung. Las fuentes correctas son las 77 páginas oficiales de Google Slides y 1,693 líneas de subtítulos hechos a mano; por eso, esta clase debe reescribirse por completo, en lugar de corregir el texto anterior con parches.

\lecturefigure{slide-001.jpg}{Tema de la clase de Nathan Lambert: Aligning Open Language Models}{official deck, p.~1}

\begin{importantbox}{La tesis central de esta clase}
El avance de la alineación abierta proviene de una cadena: un \textbf{base model más fuerte} aporta la capacidad, los \textbf{instruction/preference data} especifican la distribución de la interacción, el \textbf{fine-tuning y RLHF/DPO} modifican las probabilidades de salida, la \textbf{evaluation} proporciona la señal para iterar y la \textbf{publicación abierta} permite que la comunidad reproduzca los resultados y genere los datos de la siguiente ronda. Si cualquier eslabón de la cadena falla, ningún algoritmo por sí solo puede replicar la experiencia de ChatGPT.
\end{importantbox}

\teachervoice{Entre 00:00:34 y 00:01:03, Nathan enfatiza que no se trata de un tutorial básico, sino de volver a contar la historia del fine-tuning y la alineación posterior a ChatGPT, para que el público sepa qué conceptos siguen siendo importantes y cuáles fueron solo modas pasajeras.}

\subsection{De Shannon a ChatGPT: por qué empezar con la historia antigua}

Las siguientes diez diapositivas forman una línea de tiempo que se despliega paso a paso. Su propósito no es discutir «quién fue primero», sino mostrar que la alineación se apoya en el modelado del lenguaje, el Transformer, el escalamiento, las críticas sobre riesgos y la interacción mediante productos. Cada hito cambió lo que la comunidad abierta pudo descargar después, lo que pudo entrenar y lo que tuvo que corregir.

\lecturefigure{slide-003.jpg}{1948: Shannon aproxima el inglés con secuencias de caracteres}{official deck, p.~3}

El objetivo básico de un modelo de lenguaje autorregresivo puede escribirse como
\[
\mathcal{L}_{\mathrm{LM}}(\theta)=-\sum_{t=1}^{T}\log p_{\theta}(x_t\mid x_{<t}),
\]
donde $x_t$ es el $t$-ésimo token, $x_{<t}$ son los tokens anteriores y $\theta$ son los parámetros del modelo. Este objetivo solo exige aumentar la probabilidad de la continuation real; «ser útil», «ser inofensivo» y «obedecer las instrucciones del sistema» no son propiedades que codifique directamente.

\lecturefigure{slide-005.jpg}{2017: el Transformer convierte la attention en la estructura de cómputo principal}{official deck, p.~5}

\lecturefigure{slide-006.jpg}{2018: ELMo, GPT-1 y BERT establecen el paradigma moderno de preentrenamiento}{official deck, p.~6}

\begin{knowledgebox}{Lectura de la figura: cómo convergieron después las tres líneas}
ELMo pone el acento en las representaciones contextualizadas, BERT en la codificación bidireccional y la adaptación a tareas posteriores, y GPT en la generación autorregresiva. Los chat models actuales se desarrollaron principalmente por la línea generativa, pero incorporaron la experiencia común del aprendizaje de representaciones, el aprendizaje por transferencia y el preentrenamiento a gran escala. Lo que sustenta la línea de tiempo es la acumulación técnica, no que una sola línea estuviera destinada a imponerse desde el principio.
\end{knowledgebox}

\lecturefigure{slide-007.jpg}{2019: GPT-2, la controversia sobre la publicación de modelos y la entrada de la scaling law a la corriente principal}{official deck, p.~7}

\lecturefigure{slide-008.jpg}{La estrategia de publicación de GPT-2 convirtió a la vez la capacidad, el riesgo y la apertura en problemas}{official deck, p.~8}

\lecturefigure{slide-009.jpg}{2020: GPT-3 muestra las capacidades que trae el escalamiento y también amplía el debate sobre sesgos y daños}{official deck, p.~9}

\begin{warningbox}{El crecimiento de la capacidad y la alineación no siguen la misma curva}
Una next-token loss más baja mejora la capacidad de generación general, pero no determina por sí sola el estilo de las respuestas, los límites para negarse a responder, las citas de hechos ni las preferencias del usuario. El escalamiento amplía el «espacio de capacidades susceptibles de alinearse», pero también puede reforzar conductas no deseadas; la etapa de alineación vuelve a ponderar una distribución ya existente, en lugar de crear de la nada todas las capacidades.
\end{warningbox}

\lecturefigure{slide-010.jpg}{2021: Stochastic Parrots pone en el centro los problemas de datos, sesgos, medio ambiente y poder}{official deck, p.~10}

\lecturefigure{slide-011.jpg}{2022: ChatGPT convierte la alineación conversacional en una experiencia de producto perceptible para el público general}{official deck, p.~11}

\teachervoice{La línea de tiempo de 00:01:25 a 00:05:11 enlaza la autoregressive loss, el Transformer, el scaling, los harms y ChatGPT en un trasfondo causal. Lo que enfatiza el ponente es que no se puede explicar la alineación empezando solo por los artículos sobre RLHF, porque esta depende de más de veinte años de modelos e infraestructura previos.}

\subsection{¿Podría existir ChatGPT sin RLHF?}

Una vez que la línea de tiempo se detiene en ChatGPT, la clase plantea una pregunta más acotada: ¿es necesario el RLHF? La respuesta en clase es ``necessary, but not sufficient'': la optimización de preferencias puede cambiar de forma notable la usabilidad y el comportamiento de seguridad, pero el producto depende además del modelo preentrenado, la limpieza de datos, la prompt interface, la system policy, la retroalimentación de ingeniería y la evaluación continua.

\lecturefigure{slide-012.jpg}{Pregunta: ¿podría existir ChatGPT sin RLHF?}{official deck, p.~12}

\lecturefigure{slide-014.jpg}{Las curvas de optimización de preferencias de Anthropic muestran que el RLHF puede cambiar notablemente el comportamiento del modelo}{official deck, p.~14}

\lecturefigure{slide-015.jpg}{Llama 2 y otros modelos públicos también consideran el RLHF una etapa clave del entrenamiento}{official deck, p.~15}

\begin{knowledgebox}{Lectura de la figura: las curvas de Elo de distintas empresas no pueden compararse directamente}
En cada curva pueden diferir el prompt, los anotadores, los rivales de comparación y la escala de Elo. La figura muestra que, dentro de una misma organización, la tasa de victorias por preferencia mejora después del RLHF, pero no permite tomar los valores de la curva de Anthropic y los de la curva de Llama como si formaran una clasificación común. Antes de comparar, hay que confirmar la distribución de evaluación y el método de calibración.
\end{knowledgebox}

\teachervoice{Entre 00:05:11 y 00:06:56, Nathan usa por un lado las curvas internas de las empresas para mostrar el enorme impacto del RLHF y, por otro, advierte de inmediato que esos valores de Elo no están calibrados entre organizaciones. Por eso, la clase trata por separado «el RLHF es importante» y «cuánto pueden demostrar estas gráficas».}

\subsection{Resumen de la sección}

La alineación abierta es un sistema de moldeado de la conducta construido sobre el base model. Las diapositivas históricas muestran cómo se acumularon poco a poco las capacidades, los riesgos y las necesidades de interacción; las diapositivas sobre RLHF muestran que el entrenamiento con preferencias produce cambios de conducta notables, pero no es condición suficiente para un producto completo. El siguiente capítulo establecerá la terminología y un mapa de las secciones, para no confundir IFT, SFT, RLHF, DPO y alignment como si fueran una misma palabra.
\section{Mapa y definiciones: no confundir las tres capas Base, Instruction y Preference}

La sección anterior explicó por qué surgió la alineación; esta sección explica de qué se habla. La comunidad abierta suele usar como sinónimos ``fine-tuned'', ``chat'', ``instruct'', ``aligned'' y ``RLHF model''; el curso, en cambio, exige distinguir entre datos de entrenamiento, función de pérdida e interfaz del modelo. Solo después de separar las capas tienen sentido las comparaciones de modelos y las evaluaciones posteriores.

\subsection{El atlas del curso y los cuatro capítulos}

Esta sección revisa primero la colección de modelos y el atlas de capítulos que proporciona el ponente. Los íconos de modelos no son una genealogía completa, sino una pista reproducible: el lector puede encontrar muchos de los modelos de la clase en la Hugging Face collection. Los cuatro capítulos se desarrollan en torno a instruction, evaluation, RLHF y modern ecosystem, lo que indica que, además de los algoritmos, existe una infraestructura de datos y de medición.

\lecturefigure{slide-017.jpg}{Atlas del curso: colección de modelos, recursos de referencia y línea de tiempo de lanzamientos}{official deck, p.~17}

\lecturefigure{slide-019.jpg}{Enfoque del primer capítulo: instruction models}{official deck, p.~19}

\lecturefigure{slide-020.jpg}{El segundo capítulo agrega: expectations y evaluations}{official deck, p.~20}

\lecturefigure{slide-021.jpg}{El tercer capítulo agrega: RLHF y preference optimization}{official deck, p.~21}

\lecturefigure{slide-022.jpg}{El cuarto capítulo agrega: modern ecosystem}{official deck, p.~22}

\lecturefigure{slide-023.jpg}{Mapa completo de capítulos: instruction, evals, RLHF, ecosystem}{official deck, p.~23}

\begin{knowledgebox}{Lectura de la figura: por qué evaluation ocupa un capítulo propio}
Sin una evaluación rápida, un equipo abierto no puede saber si un cambio en los datos o en la optimización merece publicarse; sin una señal de preferencia humana a largo plazo, la evaluación automática rápida termina siendo manipulada. Evaluation no es un paso accesorio posterior al entrenamiento, sino el canal de retroalimentación que determina si el equipo puede iterar, qué modelos se vuelven visibles y qué datos se siguen recolectando.
\end{knowledgebox}

\subsection{Base model y aligned model}

Ahora dividimos los modelos en dos capas. El base model aprende principalmente una distribución general mediante next-token prediction a gran escala; el modelo que pasa por alignment, fine-tuning o preference training continúa entrenándose sobre él para que sus salidas se adapten al formato de instrucciones, a los roles de diálogo y a las preferencias. Este segundo modelo puede cambiar la experiencia de interacción con relativamente poco cómputo, pero su techo de capacidades sigue fuertemente limitado por el primero.

\lecturefigure{slide-024.jpg}{Los base models son la base de capacidades del ecosistema abierto de alineación}{official deck, p.~24}

\lecturefigure{slide-025.jpg}{Los aligned, fine-tuned y preference-trained models son la capa de adaptación de la conducta}{official deck, p.~25}

\begin{importantbox}{Diferencia de probabilidad condicional entre base y aligned}
El base model aprende $p_{\theta_0}(x_t\mid x_{<t})$; la etapa de alineación produce los parámetros $\theta$, que reordenan la probabilidad de las respuestas bajo la distribución de interacción objetivo. Conceptualmente, puede escribirse
\[
p_{\theta}(y\mid x)\propto p_{\theta_0}(y\mid x)\exp\!\left(\frac{r(x,y)}{\beta}\right),
\]
donde $r(x,y)$ es la preferencia o recompensa asignada a la respuesta y $\beta$ controla cuánto se aleja el modelo de la base/reference distribution. La fórmula expresa el objetivo de «conservar las capacidades y, al mismo tiempo, favorecer las respuestas preferidas»; no significa que todos los métodos lo implementen explícitamente de esta forma.
\end{importantbox}

\teachervoice{00:08:48--00:09:38: el ponente reconoce por iniciativa propia que la línea de tiempo omite muchas contribuciones académicas y de infraestructura, y subraya que sin base models como LLaMA/Llama 2 no existiría el ecosistema abierto de alineación posterior. Un alignment model no es un producto que exista independientemente de su base.}

\subsection{Glosario de alineación}

Esta sección descompone los términos frecuentes según «datos—objetivo—salida». Instruction fine-tuning y supervised fine-tuning suelen usar la misma autoregressive loss, pero el primero pone el énfasis en la interfaz de instrucciones; RLHF usa preferencias humanas para aprender una recompensa y optimizar la policy; DPO optimiza directamente el log-ratio policy/reference a partir de pares de preferencia. Alignment, por su parte, es un concepto de objetivo más amplio.

\lecturefigure{slide-026.jpg}{Definiciones de clase de instruction fine-tuning, SFT, RLHF, DPO y alignment}{official deck, p.~26}

\begin{knowledgebox}{Términos clave: nombre, datos y mecanismo}
\begin{tabularx}{\textwidth}{L{0.18\textwidth} L{0.24\textwidth} X}
\toprule
Término & Datos principales & Mecanismo central \\
\midrule
IFT & instruction-response pairs & Usa la LM loss para aprender el formato de instrucciones, el system prompt, los roles de varios turnos y el estilo objetivo. \\
SFT & Cualquier demostración etiquetada & Usa una pérdida supervisada para imitar la salida objetivo; no necesariamente se trata de instrucciones en lenguaje natural. \\
RLHF & preference comparisons + prompts & Aprende un reward model y luego usa métodos como PPO para maximizar la recompensa y limitar la policy drift. \\
DPO & chosen/rejected response pairs & Optimiza directamente el log-ratio de preferencia, sin ejecutar explícitamente un reward-model rollout loop. \\
Alignment & Estándares de conducta que esperan los usuarios o la organización & Puede lograrse combinando diversos datos, pérdidas, reglas y mecanismos del sistema. \\
\bottomrule
\end{tabularx}
\end{knowledgebox}

\teachervoice{00:09:53--00:11:14: Nathan define alignment de forma muy amplia: cabe en ella cualquier entrenamiento que acerque el modelo a los desires del usuario, mientras que IFT, SFT, RLHF y DPO son solo mecanismos distintos. Estas notas adoptan esa definición amplia, pero exigen precisar en cada caso los datos y la pérdida.}

\subsection{Chapter 0: la carrera inicial por reproducir ChatGPT}

Con la terminología establecida, el curso regresa a principios de 2023. En ese momento, la comunidad abierta no conocía la receta clave de los dialogue models y carecía de red-teaming, chat data y evaluaciones maduros. Numerosos equipos combinaron rápidamente base models, datos sintéticos y un fine-tuning ligero, lo que dio lugar a una carrera inicial de «un modelo por semana».

\lecturefigure{slide-027.jpg}{Chapter 0: el land grab por reproducir ChatGPT y sus problemas fundamentales}{official deck, p.~27}

\begin{warningbox}{No juzgar los primeros modelos con conocimiento retrospectivo}
El chat template, el system prompt, los datos de rechazo y la pairwise eval, que hoy parecen elementales, no tenían entonces una implementación unificada. El valor de los primeros resultados suele estar en que exponían problemas de interfaz y de datos, no en su puntuación final. El análisis histórico debe preguntar «qué desbloqueó», y no solo «si ya está obsoleto».
\end{warningbox}

\teachervoice{00:11:30--00:12:30: el ponente llama a este periodo land grab: todos publicaban modelos mientras apenas descubrían qué eran el dialogue, el red-teaming y una respuesta útil. Este contexto explica por qué los datos y la evaluación se convirtieron rápidamente en el campo de batalla principal.}

\subsection{Resumen de la sección}

Esta sección estableció la división en capas que necesita el análisis posterior: el base model determina la base de capacidades, IFT/SFT cambian la interfaz y la conducta de imitación, RLHF/DPO usan preferencias para reordenar las respuestas y alignment es un objetivo de sistema más amplio. El atlas del curso, por su parte, coloca en el mismo nivel los datos, la evaluación y los algoritmos, y nos recuerda que no debemos atribuir el avance de los modelos a un único optimizador.
\section{La primera ola de modelos instruct abiertos: la distribución de los datos importa más que el nombre}

La sección anterior definió IFT; esta sección sigue a la primera ola de modelos de 2023. Alpaca demostró que la instruction data sintética permite adaptar LLaMA con rapidez; Vicuna demostró que la distribución de prompts de chat reales es decisiva; OpenAssistant demostró que los datos humanos abiertos tienen valor a largo plazo; StableVicuna mostró que PPO no es algo que solo las empresas de código cerrado puedan ejecutar. La variable central pasó de manera continua de «si existe el algoritmo» a «si se tienen los datos correctos».

\subsection{Alpaca e instruction fine-tuning}

Esta subsección comienza con Alpaca. Su importancia no está en su calidad absoluta medida con los estándares actuales, sino en que mostró lo siguiente: un base model sólido, más una cantidad modesta de instruction-response data, basta para obtener un comportamiento chat/instruct evidente. Al leer la figura, primero distinga el conocimiento del base model del formato de respuesta aprendido mediante fine-tuning.

\lecturefigure{slide-029.jpg}{Alpaca: uno de los primeros instruction-tuned model abiertos de amplia difusión}{official deck, p.~29}

\lecturefigure{slide-030.jpg}{Los dos objetivos de IFT: adaptarse al estilo de entrada y admitir system prompt y multi-turn}{official deck, p.~30}

Si la pérdida se calcula solo sobre los assistant answer token, el objetivo de IFT puede escribirse como
\[
\mathcal{L}_{\mathrm{IFT}}=-\sum_{t\in\mathcal{A}}\log p_{\theta}(y_t\mid s,u,y_{<t}),
\]
donde $s$ es el system prompt, $u$ es la user instruction, $y$ es la assistant response y $\mathcal{A}$ son las posiciones de assistant token que requieren supervisión. El chat template codifica los roles y los límites como token especiales; una plantilla que no coincide puede hacer que el mismo modelo rinda notablemente peor.

\lecturefigure{slide-031.jpg}{IFT parte del base model y continúa el entrenamiento con instruction-response pairs}{official deck, p.~31}

\begin{importantbox}{Qué cambia principalmente IFT}
IFT suele reforzar tres tipos de capacidad: reconocer que «esto es una solicitud», organizar la respuesta según el rol y el formato, y proyectar el conocimiento ya existente hacia la salida que pide la instrucción. Puede mejorar la capacidad en la tarea, pero también puede provocar el olvido o el desplazamiento de la base capability; por ello hay que evaluar a la vez el chat behavior y las capacidades originales, y no tomar «conversa mejor» como «es más fuerte en todos los aspectos».
\end{importantbox}

\teachervoice{00:13:42--00:14:53, Nathan describe IFT como el proceso que hace que el modelo integre el chat template, el system prompt, los roles en varios turnos y el estilo de entrada. Sigue usando una autoregressive loss; la diferencia proviene sobre todo de la estructura condicional de los ejemplos de entrenamiento.}

\subsection{Self-instruct: ampliar la instruction data con modelos}

La subsección anterior explicó que el cambio de comportamiento de IFT proviene principalmente de la distribución instruction-response; esta subsección plantea la pregunta siguiente: cuando los equipos abiertos no tienen presupuesto para anotación humana a gran escala, ¿de dónde sale esa distribución? La reproducibilidad de Alpaca proviene de self-instruct: se parte de unos pocos seed prompt escritos por personas, se pide a un modelo sólido que genere más instrucciones y respuestas, y después se filtran los ejemplos duplicados, de baja calidad o con errores de formato. Al leer la figura, considere la calidad de las seed, la diversidad de la generación y las reglas de filtrado como un mismo sistema de datos; este convierte la costosa redacción humana en un pipeline de «semillas humanas + ampliación con modelos + filtrado», y también introduce el riesgo de heredar los sesgos del modelo maestro y de contraer la distribución.

\lecturefigure{slide-032.jpg}{Self-instruct: ampliar la instruction data sintética a partir de pocas seed de alta calidad}{official deck, p.~32}

Sea $D_0$ el conjunto de seed, $G$ el generador y $F$ el filtro; una ronda de ampliación puede escribirse como
\[
D_{k+1}=D_k\cup F\bigl(G(D_k)\bigr).
\]
Aquí $G(D_k)$ genera nuevos pares instruction/response y $F$ elimina duplicados y aplica las reglas de calidad. Si $G$ y $F$ solo favorecen plantillas similares, aumentar el volumen de datos no aumentará la cobertura real de tareas.

\begin{warningbox}{Los dos riesgos de ciclo cerrado de la synthetic data}
Primero, el student hereda los errores factuales, los límites de rechazo y el estilo lingüístico del teacher; segundo, la self-generation repetida comprime la diversidad y hace que el modelo sea cada vez más hábil en «preguntas al estilo de un modelo». Por ello conviene mezclar datos de usuarios reales y datos humanos, y medir la cobertura semántica de prompt/response, no solo contar el número de ejemplos.
\end{warningbox}

\teachervoice{00:14:53--00:16:10, el curso considera self-instruct un factor clave en la explosión temprana de modelos abiertos: el modelo que genera las respuestas sustituyó parte de la anotación humana y volvió accesible la instruction data. En la segunda mitad de la clase se advierte además que la distribución sintética puede ser repetitiva y frágil.}

\subsection{Vicuna y ShareGPT: los prompt de usuarios reales cambian la distribución}

Self-instruct resuelve «cómo ampliar el volumen rápidamente», pero aún no responde «si las preguntas de entrenamiento se parecen a las que harían los usuarios reales». Por eso Vicuna desplazó la atención del número de ejemplos a la prompt source: ShareGPT reunió conversaciones de ChatGPT compartidas por usuarios, lo que permitió que un modelo abierto viera por primera vez solicitudes de varios turnos, preguntas de seguimiento y estilos más cercanos a un producto real. Al revisar las dos páginas siguientes, compare a la vez la synthetic instruction y las conversaciones reales en cuanto a longitud de la tarea, estructura de turnos y riesgos de privacidad; la mejora de calidad proviene de que la distribución coincide, y no solo de tener más token.

\lecturefigure{slide-033.jpg}{Vicuna, poco después de Alpaca, incorporó una chat prompt distribution más realista}{official deck, p.~33}

\lecturefigure{slide-034.jpg}{ShareGPT: origen de los datos, valor y riesgos de privacidad}{official deck, p.~34}

\begin{knowledgebox}{Lectura de la figura: el valor de los datos y la gobernanza de los datos aparecen juntos}
ShareGPT aportó lo que los usuarios reales quieren preguntar, una señal que la synthetic instruction difícilmente puede replicar; sin embargo, una herramienta para compartir no equivale por sí misma al consentimiento del usuario para usar sus datos en entrenamiento, y el texto puede contener además información personal. Los posteriores LMSYS-Chat-1M y WildChat dieron más importancia a la limpieza y al consent, lo que muestra que la legalidad y la calidad de los datos abiertos son requisitos del mismo pipeline.
\end{knowledgebox}

\teachervoice{00:17:21--00:18:55, Nathan dice que ShareGPT es una de las pocas fuentes que permitió a los builder abiertos ver prompt de usuarios reales, y señala que los conjuntos de datos posteriores dedicaron mucho trabajo a eliminar información personal o a obtener el consentimiento antes de la recopilación.}

\subsection{Weight differences y las restricciones para publicar en abierto}

Las dos subsecciones anteriores trataron cómo se obtienen los datos; esta pasa a otra restricción del sistema que suele pasarse por alto: terminar el entrenamiento no significa que el modelo pueda entregarse de forma legal y estable. La licencia original de los pesos de LLaMA y su forma de distribución hicieron incómoda la publicación de fine-tune: los equipos solo podían publicar la weight difference respecto del base model, y los usuarios debían obtener legalmente los base weights para después combinarlos. Al leer esta línea de tiempo, distinga la capacidad del modelo, la receta de entrenamiento y el mecanismo de distribución; la weight diff no es un algoritmo de aprendizaje automático, pero determina directamente si el modelo puede instalarse, reproducirse e iterarse.

\lecturefigure{slide-035.jpg}{Los modelos instruct abiertos siguen multiplicándose: Vicuna, Koala, Dolly, entre otros}{official deck, p.~35}

\lecturefigure{slide-036.jpg}{Por qué publicar weight differences: distribuir fine-tune bajo restricciones de licencia}{official deck, p.~36}

Si los pesos base son $W_0$ y los pesos fine-tuned son $W_1$, publicar la diferencia equivale a
\[
\Delta W=W_1-W_0,\qquad W_1=W_0+\Delta W.
\]
Aquí $\Delta W$ no reduce la cantidad de información de un fine-tune completo; solo traslada al usuario la dependencia de la distribución. Es distinto del update de bajo rango de LoRA: la weight diff puede ser un tensor de rango completo y del mismo tamaño.

\lecturefigure{slide-037.jpg}{La línea de tiempo de modelos incorpora Koala y Dolly, y la forma de publicación se estandariza gradualmente}{official deck, p.~37}

\lecturefigure{slide-038.jpg}{Evolución de las puntuaciones MT-Bench de los primeros modelos instruct abiertos}{official deck, p.~38}

\begin{warningbox}{Las puntuaciones históricas solo explican el ecosistema de su momento}
El judge, los prompt y las versiones de modelo de MT-Bench pueden haberse actualizado después. Las puntuaciones de la figura sirven para ilustrar los cambios relativos entre modelos de 2023; no pueden compararse directamente con modelos de 2026 ni con configuraciones de judge distintas. Al reproducirlas, registre el commit de la evaluación, el judge model, la temperature y el prompt template.
\end{warningbox}

\teachervoice{00:19:00--00:21:05, el ponente explica que la weight diff fue una respuesta de ingeniería a las restricciones de publicación de LLaMA, y que solo los cambios posteriores de licencia facilitaron la publicación directa. La velocidad del ecosistema abierto depende de la ley y de las interfaces de distribución, no solo del código de entrenamiento.}

\subsection{OpenAssistant y StableVicuna: datos humanos y PPO temprano}

Esta subsección examina dos líneas que suelen simplificarse. OpenAssistant publicó datos de conversación generados y anotados por personas, que durante mucho tiempo fueron materia prima para modelos posteriores; StableVicuna, por su parte, completó relativamente pronto un RLHF abierto con PPO. Ambos muestran que la comunidad abierta no se limitó a IFT sintético, sino que también exploró las preferencias humanas y los pipeline de RL de alto costo.

\lecturefigure{slide-039.jpg}{OpenAssistant Conversations: uno de los primeros conjuntos de datos abiertos de instruction humanas}{official deck, p.~39}

\lecturefigure{slide-040.jpg}{StableVicuna: uno de los primeros modelos RLHF abiertos en usar PPO}{official deck, p.~40}

\begin{importantbox}{Por qué la human data es difícil de sustituir}
Los datos humanos aportan intenciones reales, lenguaje diverso, casos límite y preferencias relativas; estas señales determinan cómo el modelo ordena respuestas que son «razonables, pero no únicas». El costo proviene del reclutamiento, la capacitación, la consistencia, la privacidad y el control de calidad; por ello, lo más escaso para la comunidad abierta no suele ser el script de entrenamiento, sino datos humanos de alta calidad que puedan publicarse legalmente.
\end{importantbox}

\teachervoice{00:23:43--00:24:47, Nathan considera OpenAssistant uno de los proyectos abiertos tempranos más exitosos y presenta «se necesita más open human data» como conclusión de toda la clase; también recuerda que StableVicuna ya hizo funcionar PPO en 2023, aunque muy pocos equipos podían lograrlo.}

\subsection{Resumen de la sección}

Las variables decisivas de la primera ola de modelos instruct abiertos fueron, en orden: base weights disponibles, pipeline de instruction sintéticas, prompt de chat reales, datos humanos abiertos, licencias de publicación y un RLHF ejecutable. Los nombres de los modelos cambian sin cesar, pero la distribución de los datos y las interfaces reproducibles son una explicación más estable. La siguiente sección examinará cómo QLoRA redujo aún más la barrera de participación y por qué un fine-tuning de baja memoria no equivale automáticamente a un RLHF de bajo costo.
\section{QLoRA y Guanaco: los límites del entrenamiento con poca memoria}

Los modelos de la sección anterior todavía requieren memoria de GPU costosa e infraestructura de entrenamiento. QLoRA combina base weights cuantizados con adapters de bajo rango, de modo que modelos de 7B/13B, o incluso mayores, puedan someterse a fine-tune con menos GPU. Esto amplía de forma notable el número de participantes, pero no elimina los costos de rollout, reward modeling, RL stability y evaluación.

\subsection{La lógica de parámetros y memoria de LoRA y QLoRA}

Esta sección explica primero la low-rank adaptation. Se congelan los pesos originales $W$ y solo se entrenan dos matrices pequeñas $A,B$, con lo que un incremento de bajo rango aproxima la actualización completa. QLoRA va más allá: almacena cuantizados los base weights congelados y, durante la retropropagación, sigue calculando los gradientes a través de ellos, lo que ahorra memoria de GPU.

\lecturefigure{slide-041.jpg}{QLoRA y Guanaco: LoRA, cuantización y accesibilidad de memoria}{official deck, p.~41}

La actualización de LoRA se escribe como
\[
W'=W+\Delta W,
\qquad \Delta W=\frac{\alpha}{r}BA,
\]
donde $W\in\mathbb{R}^{d_{\mathrm{out}}\times d_{\mathrm{in}}}$ es la matriz de pesos congelada, $A\in\mathbb{R}^{r\times d_{\mathrm{in}}}$ y $B\in\mathbb{R}^{d_{\mathrm{out}}\times r}$ son matrices entrenables, $r$ es la dimensión de bajo rango y $\alpha$ es el coeficiente de escala. Cuando $r\ll d$, el número de parámetros entrenables se reduce considerablemente.

\lecturefigure{slide-042.jpg}{Comparación de la memoria de GPU que requieren full fine-tuning, LoRA y QLoRA}{official deck, p.~42}

\begin{knowledgebox}{Qué ocupa la memoria de GPU en el fine-tuning}
Un fine-tuning completo con Adam suele guardar los pesos, los gradientes y los momentos de primer orden $m$ y de segundo orden $v$ de Adam; estos optimizer states aumentan considerablemente la memoria. LoRA congela la mayoría de los pesos y solo guarda gradientes y $m,v$ para el adapter; QLoRA, además, comprime los pesos congelados a menos bits, lo que reduce todavía más la memoria residente, aunque el cómputo y el error numérico siguen exigiendo un manejo cuidadoso.
\end{knowledgebox}

La cuantización puede abstraerse como
\[
W\approx s\,(q-z),
\]
donde $q$ es un entero de pocos bits, $s$ es la scale y $z$ es el zero point. Lo esencial de QLoRA no es simplemente bajar todo el cómputo a baja precisión, sino almacenar el frozen base en representación cuantizada y, al mismo tiempo, conservar la calidad de los gradientes que necesita el entrenamiento del adapter.

\subsection{Guanaco: el valor de un método debe concretarse en un modelo publicado}

La sección anterior solo demuestra que QLoRA reduce el umbral de memoria de GPU para el fine-tuning; esta sección examina una pregunta más estricta: si la ventaja de eficiencia logra superar tres filtros (datos, optimización y evaluación) y producir al final un modelo utilizable. Que una técnica cambie de verdad el ecosistema depende de que dé lugar a un release competitivo. Guanaco entrenó modelos de varios tamaños con QLoRA y datos filtrados de OpenAssistant; al leer la figura, conviene separar primero «cuántos recursos ahorra el método» y «qué calidad alcanza el modelo publicado», y después comprobar si ambos están conectados por un mismo ciclo experimental cerrado. Guanaco muestra que un método de poca memoria no solo hace bajar la loss, sino que también puede producir un modelo conversacional de los más sólidos de su momento.

\lecturefigure{slide-043.jpg}{Guanaco: QLoRA y datos abiertos producen un modelo competitivo}{official deck, p.~43}

\begin{knowledgebox}{Lectura de la figura: los tres niveles de evidencia de un artículo de métodos}
El primer nivel es si disminuyen la memoria y aumenta el rendimiento; el segundo, si el entrenamiento converge de forma estable; el tercero, si el modelo final es competitivo en evaluaciones humanas o confiables. Solo cuando aparece el tercer nivel la comunidad sabe que el método no sacrificó calidad esencial a cambio de eficiencia. La importancia de Guanaco está en haber conectado los tres niveles.
\end{knowledgebox}

\teachervoice{00:24:56--00:27:09, Nathan describe QLoRA como la técnica que abrió la participación a más actores y compara una A100 de 80GB con la memoria de las GPU de consumo. El ponente subraya también que fueron Guanaco y la versión filtrada de los datos de OpenAssistant lo que convirtió el método en la publicación de un modelo confiable.}

\subsection{Por qué los LoRA methods no tuvieron éxito automático en RL}

Una vez que el IFT con poca memoria se volvió fácil de reproducir, muchos equipos intentaron aplicar LoRA/QLoRA directamente a RLHF. La dificultad es que RL requiere policy rollout, reference policy, reward model, advantage estimation y múltiples actualizaciones, de modo que la memoria no es el único cuello de botella; la parametrización con adapters también modifica la geometría de la optimización.

\lecturefigure{slide-045.jpg}{¿Son adecuados los LoRA methods para RL? Mucha exploración, pocos modelos que trascendieron}{official deck, p.~45}

\begin{warningbox}{Que baje la loss no demuestra que RLHF haya funcionado}
Un RL pipeline puede hacer bajar el objetivo de entrenamiento y, aun así, producir reward hacking, deriva de KL, colapso de modos o degradación en las evaluaciones. Se necesitan un release del modelo final, evaluaciones independientes y muestras de comportamiento. En la sesión de preguntas, Nathan da una regla empírica: antes de dedicarse a fondo a un método, esperar a que produzca un modelo competitivo de magnitud comparable.
\end{warningbox}

\teachervoice{00:27:28--00:28:45, el ponente recuerda que en ese momento la loss de LoRA-RL bajaba y los experimentos entusiasmaban, pero hubo muy pocos modelos que realmente causaran un «splash». A partir de 01:05:44 añade que él espera a que los métodos de este tipo produzcan un release competitivo antes de profundizar en ellos.}

\begin{lstlisting}[caption={Compuertas de validación: de un método eficiente a la publicación de un modelo confiable}]
check_memory_reduction()
check_training_stability()
check_base_capability_retention()
run_fast_automatic_evals()
run_human_or_arena_evaluation()
release_model_recipe_and_data_provenance()
\end{lstlisting}

\subsection{Resumen de la sección}

Lo que resuelve QLoRA es el umbral de participación: reduce los parámetros entrenables y el almacenamiento de los frozen weights, de modo que más equipos puedan hacer IFT. No resuelve automáticamente los datos, el rollout, la estabilidad ni la evaluación de RLHF. El criterio de evidencia del curso es el ciclo cerrado «método—modelo—evaluación», no solo las curvas de recursos o la loss de entrenamiento.
\section{Safety Backlash y periodo de transición: el propio límite de las negativas a responder es una decisión de alineación}

Además de los algoritmos y los datos, la alineación abierta enfrenta conflictos de valores. La política de seguridad de Llama 2 Chat provocó el rechazo de una parte de los usuarios, y después aparecieron los fine-tunes «uncensored»; en el mismo periodo, muchos modelos como UltraChat, OpenChat, XwinLM y OpenHermes iteraron rápidamente con distintos datos y recipes. El curso considera esta etapa una franja de transición en la que las expectations todavía no se estabilizaban.

\subsection{Llama 2 safety backlash}

El ciclo «método—modelo—evaluación» de la sección anterior supone que existe consenso sobre qué es un buen comportamiento, pero el backlash de seguridad muestra precisamente que esa premisa no se cumple. Esta sección no busca responder si el modelo necesita restricciones de seguridad, sino quién define el límite de esas restricciones y quién asume el costo de los rechazos erróneos. El modelo debe trazar una frontera entre las solicitudes peligrosas, las preguntas sensibles pero legítimas, el humor y la ficción; si los datos premian en exceso la negativa a responder, el modelo también se retraerá ante solicitudes inofensivas. Al leer la figura, no hay que interpretar el rechazo de la comunidad como «la seguridad no sirve», sino advertir que el comportamiento de seguridad debe evaluarse junto con el escenario del usuario, las políticas, la controlabilidad y el costo de los rechazos erróneos.

\lecturefigure{slide-046.jpg}{El safety backlash de Llama 2 Chat: qué tan seguro debe ser un chat model}{official deck, p.~46}

\teachervoice{00:28:55--00:29:52, Nathan considera el backlash de Llama 2 un defining moment: el estilo de negativas del modelo hizo que la comunidad se diera cuenta de que «safe» no es un valor numérico único, sino una divergencia entre usuarios y desarrolladores sobre la controlabilidad.}

\subsection{Objetivos y riesgos de los «Uncensored» models}

Después del backlash de Llama 2, la comunidad propuso pronto una dirección opuesta, practicable pero burda: retirar de los datos de entrenamiento las negativas explícitas. Esta decisión pone el «control del usuario» en una posición más alta, pero no aporta por sí misma un nuevo mecanismo para juzgar el daño. Algunos fine-tunes buscan no negarse nunca mediante la eliminación de ejemplos que contienen ``Sorry'', ``as a language model'' y expresiones similares; al leer la figura, lo que debe compararse es el refusal style frente a la harmful compliance real, no solo contar frases de disculpa. Estos modelos demuestran que los datos de entrenamiento pueden cambiar rápidamente la refusal distribution, y también revelan que «quitar las negativas superficiales» no equivale a construir un mejor razonamiento de seguridad; el modelo podría aprender únicamente a recurrir menos a las plantillas de seguridad.

\lecturefigure{slide-047.jpg}{«Uncensored» models: eliminar ejemplos de negativas para aumentar el control del usuario}{official deck, p.~47}

\begin{warningbox}{La tasa de negativas no es una safety metric completa}
Hay que distinguir, como mínimo, la harmful compliance, la benign refusal, los errores factuales, la filtración de datos privados y la reversibilidad. Una tasa baja de negativas puede mejorar la utilidad en tareas legítimas, pero también puede aumentar la ayuda peligrosa; una tasa alta puede ser más segura, pero también puede ser solo una evasión por plantilla. La evaluación debe estratificarse por categoría de riesgo y por intent del usuario.
\end{warningbox}

\teachervoice{00:29:52--00:30:54, el ponente explica que los modelos uncensored cambian el behavior al retirar los ejemplos de entrenamiento con negativas, una línea que continúa hasta hoy. La clase no lo presenta como una respuesta de seguridad, sino como evidencia de las decisiones de valores de la comunidad.}

\subsection{Modelos del periodo de transición: las recetas se difunden, pero falta una señal unificada}

A medida que se difundieron la instruction data y el código de entrenamiento, el número de modelos creció rápidamente, pero resultaba difícil determinar qué recipe era realmente mejor. UltraChat, OpenChat, XwinLM, OpenHermes y otros combinaron datos sintéticos, conversaciones reales, distintos base models y optimización de preferencias; las diferencias entre modelos se fueron convirtiendo en un problema de evaluación.

\lecturefigure{slide-048.jpg}{UltraChat, OpenChat, XwinLM, OpenHermes y otros modelos del periodo de transición}{official deck, p.~48}

\begin{knowledgebox}{El verdadero producto del periodo de transición}
Lo que dejó esta etapa no fueron solo modelos, sino conjuntos de datos, chat templates, marcos de entrenamiento y puntos de entrada para la evaluación. Muchos modelos quedaron obsoletos pronto, pero la infraestructura se convirtió en componente de Zephyr, Tulu y las post-training recipes modernas. Para juzgar una contribución histórica hay que mirar los artifacts reutilizables, no solo la posición en la clasificación el día del lanzamiento.
\end{knowledgebox}

\subsection{Resumen de la sección}

El safety backlash muestra que la alignment es una decisión sobre la distribución de salidas y el límite de las negativas; los modelos del periodo de transición muestran que, cuando las recetas se difunden rápidamente, lo que más le falta a la comunidad es retroalimentación confiable. La siguiente sección presenta cuatro herramientas de evaluación y compara si las preferencias humanas, LLM-as-a-judge y los benchmarks estáticos se adecúan mejor a decisiones de desarrollo o de lanzamiento.
\section{Evaluation Infrastructure: sin retroalimentación no hay iteración de la alineación abierta}

La sección anterior presentó una gran cantidad de modelos, pero sin un sistema de referencia común. Entre mayo y julio de 2023 aparecieron, en cuestión de semanas, ChatBotArena, AlpacaEval, MT-Bench y Open LLM Leaderboard. La causa no fue que la teoría de la evaluación madurara de repente, sino que los equipos abiertos no podían, como las grandes empresas, contratar personas de forma continua para comparar respuestas. Por eso la evaluación se convirtió en la retroalimentación sustituta del loop de entrenamiento.

\subsection{Por qué surgieron al mismo tiempo cuatro sistemas de evaluación}

Esta sección revisa primero la cronología y la división de funciones. Arena recopila pairwise preference de personas reales: su retroalimentación es lenta, pero tiene alta validez externa. AlpacaEval y MT-Bench usan un LLM judge potente para acelerar la iteración. Leaderboard ejecuta automáticamente tareas estáticas y es el más fácil de integrar en un pipeline de ingeniería. Los cuatro optimizan objetivos distintos y ninguna puntuación única puede sustituirlos.

\lecturefigure{slide-049.jpg}{En 2023, cuatro de las principales aligned-model evaluations abiertas aparecieron casi al mismo tiempo}{official deck, p.~49}

\teachervoice{00:32:14--00:32:48, Nathan explica que estas herramientas surgieron en poco tiempo porque los builders abiertos desperately needed signal, pero no podían pagar de forma continua el costo de las comparaciones humanas como lo hacen las empresas de código cerrado.}

\subsection{ChatBotArena: lenta, pero cercana a la preferencia real de los usuarios}

La subsección anterior colocó los cuatro sistemas de evaluación en un mismo mapa; esta examina primero el que más se acerca a la retroalimentación real de un producto. Arena muestra dos modelos anónimos lado a lado al usuario, recopila victorias, derrotas o empates y después estima su fuerza relativa mediante pairwise ranking. Al leer el diagrama de flujo, conviene seguir primero de dónde viene el prompt, cómo entran los resultados de la comparación en el ranking y cuántos enfrentamientos necesita un modelo nuevo para estabilizarse. Su ventaja es que los prompts vienen de usuarios reales y el judge es una persona; su debilidad es que tanto el muestreo como la convergencia son lentos, de modo que la retroalimentación para el desarrollo puede tardar varios días.

\lecturefigure{slide-050.jpg}{ChatBotArena: comparación anónima entre dos modelos y recopilación de la preferencia real de los usuarios}{official deck, p.~50}

La probabilidad esperada de victoria en Elo puede escribirse como
\[
P(A\succ B)=\frac{1}{1+10^{(R_B-R_A)/400}},
\]
donde $R_A,R_B$ son las calificaciones de los modelos y $P(A\succ B)$ es la probabilidad de que A venza a B. La calificación es una magnitud relativa que depende del grafo de oponentes, de la distribución de usuarios y de la ventana temporal; no conviene comparar directamente el Elo absoluto entre distintas versiones de Arena.

\teachervoice{00:33:14--00:33:56, el ponente considera que Arena es muy importante para los lanzamientos y la estrategia de las empresas, pero demasiado lenta para la ingeniería de entrenamiento: el equipo necesita saber si el modelo mejoró antes de decidir el lanzamiento, por lo que debe complementarse con evaluaciones automáticas más rápidas.}

\subsection{AlpacaEval: preference proxy rápido y judge bias}

AlpacaEval usa prompts fijos, compara la candidate response con la baseline response y después pide a un LLM potente que decida cuál gana. Más muestras producen barras de error más pequeñas; es rápido de ejecutar y fácil de integrar. Sin embargo, la baseline, el judge, la response length y el style modifican la tasa de victorias.

\lecturefigure{slide-051.jpg}{AlpacaEval: LLM preference judging entre la candidate response y la baseline response}{official deck, p.~51}

\lecturefigure{slide-052.jpg}{Ventajas y limitaciones de AlpacaEval: muchas muestras y poco error, pero el judge tiene sesgos}{official deck, p.~52}

Si en $n$ comparaciones independientes la proporción de victorias de la candidate es $\hat p$, el error estándar aproximado es
\[
\operatorname{SE}(\hat p)=\sqrt{\frac{\hat p(1-\hat p)}{n}}.
\]
donde $n$ es el número de comparaciones y $\hat p$ es la win rate observada. Aumentar las muestras reduce el error aleatorio, pero no elimina el sesgo sistemático; por ejemplo, que el judge prefiera respuestas más largas, más corteses o parecidas a su propio estilo.

\lecturefigure{slide-053.jpg}{AlpacaEval 2: cambio a una baseline de GPT-4 e incorporación de length control}{official deck, p.~53}

\begin{warningbox}{Cuatro fuentes de sesgo en LLM-as-a-judge}
\begin{enumerate}
  \item \textbf{Position bias}: la respuesta que aparece primero o al final gana con más facilidad;
  \item \textbf{Length/style bias}: las respuestas extensas y con formato vistoso reciben una preferencia adicional;
  \item \textbf{Self preference}: el judge prefiere salidas parecidas a su propia distribución;
  \item \textbf{Contamination}: el judge puede haber visto el benchmark, la candidate o la reference.
\end{enumerate}
El intercambio aleatorio de posiciones, el control de longitud, el uso de varios judges y la revisión humana por muestreo solo atenúan estos sesgos; no los eliminan por completo.
\end{warningbox}

\teachervoice{00:33:56--00:36:16, Nathan considera AlpacaEval una señal de ingeniería rápida, pero se muestra cauteloso sobre el significado concreto de las puntuaciones de AlpacaEval 2. Cuanto más fácil de optimizar es una evaluación, más necesario es comprobar que el modelo no solo aprendió a complacer al judge.}

\subsection{MT-Bench: cobertura multiturno y fragilidad por muestra pequeña}

AlpacaEval usa pairwise preference para evitar la calibración de puntuaciones absolutas; MT-Bench, en cambio, opta deliberadamente por calificaciones escalares multiturno y por categorías, para observar si el modelo mantiene la calidad después de preguntas de seguimiento. Por eso esta subsección compara dos fuentes de error: la preferencia relativa del pairwise judge y la deriva de una escala de 1 a 10. MT-Bench pide a un LLM judge que califique respuestas de dos turnos, con preguntas que cubren categorías como redacción, razonamiento, matemáticas y programación. Al leer la figura, conviene comprobar primero si el segundo turno depende realmente del primero y después si el promedio por categoría oculta fallos individuales. La estructura multiturno se acerca más al chat que la preference de un solo turno, pero el número de preguntas es pequeño, la escala de calificación del judge es inestable y los equipos pueden ajustar fácilmente sus modelos a las preguntas públicas.

\lecturefigure{slide-054.jpg}{MT-Bench: un modelo potente asigna calificaciones escalares a respuestas multiturno}{official deck, p.~54}

\lecturefigure{slide-055.jpg}{Ventajas, desventajas y riesgo de sobreoptimización de MT-Bench}{official deck, p.~55}

\begin{knowledgebox}{Lectura de la figura: las puntuaciones escalares son más difíciles de calibrar que las pairwise}
Asignar un 7 a una respuesta requiere un criterio absoluto estable, mientras que juzgar si A es mejor que B solo requiere una comparación local. Distintos judges o distintas prompt wording alteran la escala de 1 a 10. Al reportar MT-Bench deben conservarse los resultados por pregunta, la versión del judge y la distribución por categorías, no solo el promedio.
\end{knowledgebox}

\teachervoice{00:36:16--00:37:44, el curso señala que las preguntas de dos turnos de MT-Bench son útiles, pero son pocas, el judge tiene sesgos y, una vez publicadas, pueden optimizarse directamente. Un modelo con puntuación alta puede simplemente ajustarse mejor a esas 80 preguntas.}

\subsection{Open LLM Leaderboard: cómo una herramienta de ingeniería cae bajo la ley de Goodhart}

Leaderboard ejecuta automáticamente varios benchmarks estáticos; en un principio era una herramienta de ingeniería para comparar con facilidad modelos base/fine-tuned. A medida que la clasificación empezó a influir en las descargas y en la reputación, los datos de entrenamiento comenzaron a incluir contenido con el estilo de los benchmarks e incluso preguntas de prueba, y las puntuaciones estáticas fueron perdiendo la capacidad de distinguir la chat quality real.

\lecturefigure{slide-056.jpg}{Open LLM Leaderboard: pipeline automático de benchmarks y clasificación pública}{official deck, p.~56}

\lecturefigure{slide-057.jpg}{Combinación de evaluaciones: Arena es lenta y realista; las eval automáticas son rápidas y se prestan a la manipulación}{official deck, p.~57}

\begin{importantbox}{La forma que toma la ley de Goodhart en la evaluación de modelos}
Cuando una metric pasa de «objetivo de medición» a «objetivo de entrenamiento», el modelo aprende la parte aprovechable de la metric. La solución no es buscar un único benchmark que nunca pueda manipularse, sino construir una combinación de evaluaciones: prompts ocultos o dinámicos, preferencia real de los usuarios, estratificación entre capacidades y seguridad, auditoría humana y reemplazo periódico de las tareas ya saturadas.
\end{importantbox}

\teachervoice{00:37:44--00:39:45, Nathan explica que Leaderboard fue en un principio una herramienta de ingeniería, pero la contaminación y el gaming redujeron su valor; su recomendación de combinación es usar Arena como señal pública de largo plazo y AlpacaEval/MT-Bench como señales rápidas de desarrollo, en lugar de buscar una puntuación universal.}

\begin{lstlisting}[caption={Estructura mínima de registro de un portfolio de evaluación}]
evaluation = {
    "prompt_distribution": "real_users | fixed_set | hidden_set",
    "judge": "human | model_name_and_version",
    "baseline": "model_or_none",
    "aggregation": "elo | win_rate | scalar_mean",
    "latency": "minutes | days",
    "known_biases": ["length", "position", "contamination"]
}
\end{lstlisting}

\subsection{Resumen de la sección}

Los cuatro sistemas de evaluación sirven a escalas de tiempo distintas: Arena respalda los lanzamientos y la comparación pública; AlpacaEval/MT-Bench respaldan la iteración rápida, y Leaderboard respalda la regresión automática de capacidades. Cada herramienta parte de supuestos sobre el prompt, el judge, la baseline y la aggregation. La madurez de la alineación abierta no consiste en encontrar una puntuación perfecta, sino en saber qué puede responder cada puntuación y qué no.
\section{Objetivo de RLHF, reward model y DPO: fórmulas parecidas, procesos de optimización distintos}

Con la retroalimentación de la evaluación disponible, esta sección pasa al entrenamiento con preferencias. RLHF primero entrena un reward model con comparaciones humanas y después optimiza la policy; DPO convierte el mismo tipo de señal de preferencia en una pérdida directa de tipo clasificación. Ambos buscan aumentar la probabilidad de la chosen response frente a la rejected response, pero difieren en el muestreo de datos, el rollout en línea, la estabilidad y el costo de cómputo.

\subsection{RLHF objective: recompensa y restricción de KL}

Esta subsección empieza por el objetivo estándar. Si solo se maximiza la learned reward, la policy puede encontrar fallas del reward model y alejarse del base behavior; por eso se añade la KL divergence respecto al reference model, que limita la deriva de la distribución. PPO es el optimizador habitual, pero el objetivo en sí no equivale a PPO.

\lecturefigure{slide-059.jpg}{Notación del RLHF objective: policy, reference, prompt, response y reward}{official deck, p.~59}

\lecturefigure{slide-060.jpg}{Primera parte del RLHF objective: maximizar la puntuación del reward model}{official deck, p.~60}

\lecturefigure{slide-061.jpg}{Segunda parte del RLHF objective: usar KL para impedir que la policy se aleje de la reference}{official deck, p.~61}

El objetivo se escribe como
\[
\max_{\theta}\;\mathbb{E}_{x\sim D,\,y\sim\pi_{\theta}(\cdot\mid x)}
\left[r_{\phi}(x,y)-\beta\log\frac{\pi_{\theta}(y\mid x)}{\pi_{\mathrm{ref}}(y\mid x)}\right].
\]
donde $x$ es el prompt, $y$ es la respuesta muestreada por la policy $\pi_{\theta}$, $r_{\phi}$ es el reward model, $\pi_{\mathrm{ref}}$ es el reference model y $\beta$ controla la KL penalty. Si $\beta$ es demasiado pequeño, aparecen la deriva y el reward hacking; si es demasiado grande, resulta difícil cambiar el comportamiento.

\teachervoice{00:40:00--00:41:55, el ponente descompone la fórmula en dos fuerzas, «perseguir la reward» y «no alejarse demasiado de la base», y enfatiza que la reference constraint no es un adorno, sino lo que impide que la learned reward se explote en exceso.}

\subsection{Preference modeling: por qué se usa la pairwise comparison}

A una persona le cuesta asignar de forma estable una puntuación absoluta de 0--10 a una respuesta abierta, pero le resulta más fácil juzgar cuál es mejor entre A y B. Por eso el reward model suele usar chosen/rejected pairs y aprende un orden relativo. La calidad de los datos depende del prompt, de la forma de generar los candidatos, de la guía de anotación y del annotator agreement.

\lecturefigure{slide-062.jpg}{Preference reward modeling: la comparación relativa supera a la supervisión directa de puntuaciones absolutas}{official deck, p.~62}

El modelo de Bradley--Terry se escribe como
\[
P(y_w\succ y_l\mid x)=\sigma\!\left(r_{\phi}(x,y_w)-r_{\phi}(x,y_l)\right),
\]
donde $y_w$ es la preferred/chosen response, $y_l$ es la rejected response, $r_{\phi}$ es una reward escalar y $\sigma$ es la sigmoide. El entrenamiento minimiza la log-probabilidad negativa de que la chosen gane.

\begin{warningbox}{Los preference data no equivalen a una verdad objetiva}
Las anotaciones reflejan una población concreta, unas políticas, una interfaz y un conjunto de candidatos. Si ambos candidatos son malos, la pairwise label de todos modos elige uno; si la diferencia entre candidatos es demasiado evidente, el reward model no aprende las fronteras finas. Conviene conservar tie/skip, la información de los anotadores y el disagreement, en lugar de guardar solo etiquetas binarias.
\end{warningbox}

\teachervoice{00:42:00--00:43:44, Nathan explica que supervisar directamente una scalar reward no da buenos resultados, mientras que la pairwise preference es más natural. El punto del curso no es que las comparaciones humanas carezcan de sesgo, sino que el juicio relativo suele ser más estable que la puntuación absoluta.}

\subsection{De la «gradient ascent directa» a DPO}

Con un pairwise reward objective, cabe preguntarse si es indispensable entrenar explícitamente un reward model y después ejecutar RL. DPO aprovecha el óptimo KL-regularized, escribe la reward implícita como el log-ratio entre policy y reference, y aplica directamente una logistic classification a cada chosen/rejected pair.

\lecturefigure{slide-063.jpg}{Pregunta: ¿se puede aplicar gradient ascent directamente al preference objective?}{official deck, p.~63}

\lecturefigure{slide-064.jpg}{Respuesta: Direct Preference Optimization reescribe el objetivo como una pérdida de preferencia directa}{official deck, p.~64}

La DPO loss es
\[
\mathcal{L}_{\mathrm{DPO}}(\theta)=-\mathbb{E}\log\sigma\!\left(
\beta\left[
\log\frac{\pi_{\theta}(y_w\mid x)}{\pi_{\mathrm{ref}}(y_w\mid x)}-
\log\frac{\pi_{\theta}(y_l\mid x)}{\pi_{\mathrm{ref}}(y_l\mid x)}
\right]\right).
\]
donde $\pi_{\theta}$ es la policy que se entrena, $\pi_{\mathrm{ref}}$ es la reference, $y_w,y_l$ son la chosen y la rejected, y $\beta$ controla la intensidad con que se permite alejarse de la reference. DPO aumenta el log-ratio relativo de la chosen y, al mismo tiempo, reduce el log-ratio relativo de la rejected.

\teachervoice{00:43:44--00:46:20, el ponente resume el atractivo de DPO en una implementación mínima que no requiere un RL rollout explícito; pero enseguida enfatiza que DPO y PPO son optimizers muy distintos: que la derivación de sus objetivos esté conectada no significa que su dinámica de entrenamiento sea la misma.}

\subsection{DPO core facts y diferencias con PPO}

Esta subsección separa la comodidad de implementación de los límites de capacidad. DPO solo necesita preference pairs fuera de línea y su entrenamiento se parece a una supervised classification; PPO muestrea en línea a partir de la policy actual, usa la reward para estimar la advantage y realiza actualizaciones de varios pasos mediante clipping/regularization. PPO es más complejo y más costoso, y también puede extraer más de la reward signal.

\lecturefigure{slide-065.jpg}{DPO core facts: implementación sencilla, adopción rápida, pero con supuestos}{official deck, p.~65}

\lecturefigure{slide-066.jpg}{DPO y PPO/REINFORCE son optimizadores y ciclos de datos distintos}{official deck, p.~66}

\begin{knowledgebox}{Términos clave: objective y optimizer}
\begin{tabularx}{\textwidth}{L{0.18\textwidth} X}
\toprule
Concepto & Significado \\
\midrule
Preference objective & Busca que la chosen sea más probable que la rejected y, a la vez, limita cuánto se aleja de la reference. \\
DPO & Optimiza directamente el log-ratio policy/reference con pairs fijos fuera de línea. \\
PPO & Hace rollout desde la policy actual y realiza actualizaciones on-policy restringidas con base en la reward/advantage. \\
REINFORCE & Policy-gradient estimator clásico que multiplica el retorno muestreado por el score-function gradient. \\
\bottomrule
\end{tabularx}
Un mismo objetivo de alto nivel puede aproximarse con distintos optimizers; la estabilidad, la eficiencia de muestras y el comportamiento final no tienen por qué coincidir.
\end{knowledgebox}

\teachervoice{En la Q\&A, 01:09:02--01:09:40, Nathan ofrece un juicio basado en la experiencia: si PPO se logra ajustar bien, a menudo extrae un poco más de rendimiento de los mismos datos; pero la ventaja se ubica en fine margins y su costo es una gran cantidad de modelos e iteraciones de entrenamiento.}

\subsection{Resumen de la sección}

El núcleo de RLHF es la learned reward junto con la reference constraint; la pairwise preference es la forma de datos habitual del reward modeling; DPO reescribe la reward implícita como una pérdida directa de clasificación de preferencias. La conexión matemática no borra las diferencias entre optimizers. Para evaluar DPO/PPO hay que considerar el release final, los datos, el cómputo, la estabilidad y las preferencias humanas, no solo la extensión de la implementación.
\section{Zephyr, Tulu 2, SteerLM y Starling: no hay un ganador único}

La sección anterior presentó las fórmulas; esta sección examina la evidencia de los modelos. Zephyr demostró que DPO puede producir un modelo abierto que trasciende el ámbito especializado, y Tulu 2 extendió DPO a 70B; SteerLM y Starling, por su parte, demostraron que la vía de PPO/RL todavía puede ser más fuerte. La conclusión histórica no es «DPO sustituye a PPO», sino que la preference optimization empezó a producir releases públicas comparables.

\subsection{Zephyr y Tulu 2: DPO pasa del artículo a la release}

Este apartado revisa primero dos implementaciones decisivas. Zephyr combina UltraChat y UltraFeedback y le da a DPO una recipe clara y reproducible; Tulu 2 usa los datos y la infraestructura de entrenamiento de AI2 para extender el método a 70B. La publicación de los modelos permite a la comunidad comprobar si el método conserva su valor en distintas escalas y con distintos datos.

\lecturefigure{slide-067.jpg}{Zephyr $\beta$: uno de los primeros modelos abiertos que recibió amplia atención gracias a DPO}{official deck, p.~67}

\lecturefigure{slide-068.jpg}{Tulu 2: extensión de DPO a la escala de 70B}{official deck, p.~68}

\begin{knowledgebox}{Lectura de la figura: las partes transferibles de la recipe}
Al comparar releases hay que descomponerlas en base model, IFT data, preference data, reference policy, $\beta$, sequence length y eval. Si Zephyr y Tulu cambiaron varios componentes a la vez, no se puede atribuir toda la mejora a DPO; su aportación consiste en ofrecer un artifact completo para que los equipos posteriores puedan hacer ablaciones componente por componente.
\end{knowledgebox}

\teachervoice{00:46:33--00:48:30, Nathan describe a Zephyr como el primer modelo con el que DPO realmente «make a splash», y considera a Tulu 2 un hito de escalamiento. La credibilidad de un método ante la comunidad proviene de la release de modelos, no solo de la derivación en un artículo.}

\subsection{SteerLM y Starling: PPO/RL sigue siendo competitivo}

Que DPO sea sencillo no significa que el RL on-policy haya dejado de funcionar. SteerLM controla la salida mediante condicionamiento por atributos y señales de preferencia, y Starling ejecuta PPO sobre datos de preferencia al estilo de Arena, lo que muestra que un pipeline más complejo todavía puede obtener ventajas. El curso coloca deliberadamente estos modelos después del éxito de DPO para impedir el relato de que «un método nuevo lo unifica todo».

\lecturefigure{slide-069.jpg}{SteerLM y Starling: los métodos PPO/RL todavía pueden superar a los modelos DPO}{official deck, p.~69}

\begin{warningbox}{Las clasificaciones de algoritmos suelen disfrazar una ventaja de datos como ventaja del optimizador}
Si el modelo PPO usa preference data más reciente, más grande o más realista, mientras que la baseline DPO usa datos antiguos, el resultado no se puede atribuir al optimizer. Una comparación confiable requiere compartir el checkpoint base/IFT, compartir los pairs o el reward model, tener un presupuesto de cómputo similar y usar el mismo eval portfolio.
\end{warningbox}

\teachervoice{00:48:45--00:49:35, el expositor usa SteerLM y Starling para señalar que «todavía hay muchos modelos que muestran que PPO/RL supera a DPO». Su conclusión es que la evidencia es mixta, no tomar partido por un bando.}

\subsection{Resumen de la sección}

Zephyr y Tulu 2 demostraron la utilidad práctica de DPO y su escalamiento, mientras que SteerLM/Starling mantuvieron la competitividad de PPO. En la alignment abierta, la unidad experimental adecuada no es el nombre del algoritmo, sino la recipe completa. La siguiente sección entra en el ecosistema de 2024 y analiza la llegada de más participantes, Llama 3, la brecha entre modelos abiertos y cerrados, y el verdadero cuello de botella: la preference data de alta calidad.
\section{El ecosistema moderno de alineación abierta: más participantes, pero todavía faltan datos}

Las secciones anteriores siguieron los modelos en orden cronológico; esta sección toma abril de 2024 como una instantánea de la clase. Los participantes ya no son solo universidades, organizaciones sin fines de lucro y la comunidad de Hugging Face, sino también numerosas empresas; las recipe de los modelos son más complejas, la evaluación es más rápida y la brecha entre modelos abiertos y cerrados se ha reducido en algunas tareas. Sin embargo, el ponente considera que la limitación más persistente sigue siendo los datos, en particular la preference data humana legal, abierta y diversa.

\subsection{Más modelos y más roles en la generación de datos}

Esta subsección muestra la diversidad del ecosistema moderno. Distintos equipos se especializan en instruction rewriting, preference generation, filtrado, chat data, reward model y frameworks de entrenamiento. La alignment abierta ya no es un trabajo de extremo a extremo de un solo equipo, sino una artifact supply chain.

\lecturefigure{slide-071.jpg}{Diversificación de los roles de generación de datos, entrenamiento de modelos y comunidad en el ecosistema moderno}{official deck, p.~71}

\begin{knowledgebox}{Artifact supply chain}
Un aligned model abierto puede depender de: base weights de terceros, instructions generadas por otro modelo, chat logs públicos, preference data de la comunidad, un trainer de código abierto, un LLM judge cerrado y un leaderboard público. Cada dependencia introduce riesgos de licencia, contaminación, estilo y reproducibilidad, por lo que la model card debe registrar el lineage completo.
\end{knowledgebox}

\subsection{Llama 3: sobre todo una scale story, pero también una post-training signal}

Llama 3 se publicó en las fechas de la clase, y el ponente lo considera «más sobre scaling» que un único breakthrough de alignment. Meta reporta el uso de métodos en múltiples etapas, como IFT, rejection sampling, DPO y PPO; un pipeline tan complejo muestra que la industria está dispuesta a apilar métodos para obtener mejoras pequeñas, y también que todavía no se ha encontrado un algoritmo simple y unificado.

\lecturefigure{slide-072.jpg}{Llama 3: instantánea de la clase sobre el preentrenamiento a gran escala y el post-training en múltiples etapas}{official deck, p.~72}

Conceptualmente, la rejection sampling puede escribirse así: se muestrean $K$ candidatos de la policy
\[
y_1,\ldots,y_K\sim\pi(\cdot\mid x),
\qquad y^{\star}=\arg\max_{y_i} r(x,y_i),
\]
donde $r$ puede ser un reward model, un conjunto de reglas o un judge. Reincorporar $y^{\star}$ al IFT puede mejorar la inicialización, pero también fija en los nuevos datos los sesgos del reward o del judge.

\teachervoice{01:02:44--01:03:22, a Nathan le parece compleja la receta de Llama 3, que combina IFT, rejection sampling, DPO y PPO; conjetura que estas etapas van desplazando gradualmente las capacidades y sirven de inicialización para la etapa siguiente, y que en el futuro podrían destilarse en un algoritmo más simple.}

\subsection{Current directions: la alignment se está convirtiendo en un data problem}

Esta subsección pasa de los modelos a las preguntas de investigación. La comunidad abierta cuenta con pocos preference datasets de uso extendido, como Anthropic HH, UltraFeedback y Nectar; su reutilización repetida provoca sobreajuste, homogeneización y contamination. En el futuro se necesitarán más datos de interacción reales con personas, multilingües, de múltiples dominios y de largo plazo.

\lecturefigure{slide-073.jpg}{Current directions: la frontera de investigación en la alineación de modelos de lenguaje abiertos}{official deck, p.~73}

\lecturefigure{slide-074.jpg}{La brecha de capacidades entre aligned models abiertos y cerrados varía según la tarea}{official deck, p.~74}

\begin{warningbox}{«Open catches closed» debe especificar la tarea y la fecha}
Que algunos modelos abiertos se acerquen a los cerrados en benchmarks públicos no significa que los hayan alcanzado en uso de herramientas, contexto largo, confiabilidad, seguridad, latencia y capacidad multilingüe. La figura es un snapshot de abril de 2024; las conclusiones cambian al actualizarse las versiones de los modelos y los judges. Toda comparación debe indicar el exact model, la date y el evaluation protocol.
\end{warningbox}

\lecturefigure{slide-075.jpg}{Primera dirección actual: la preference data es muy insuficiente y la calidad de los datos limita los experimentos}{official deck, p.~75}

\begin{importantbox}{Matriz de cobertura de un preference dataset}
Los datos de alta calidad deben cubrir, como mínimo: dominio de la tarea, idioma, grupo de usuarios, nivel de riesgo, longitud de la respuesta, con herramientas o sin ellas, estado de múltiples turnos y disagreement. Aumentar solo el número de muestras de una misma synthetic template reduce la varianza, pero no amplía los límites del comportamiento.
\end{importantbox}

\teachervoice{01:00:00--01:03:48, el ponente señala la escasez de preference data como el cuello de botella central actual, y advierte que la synthetic data suele ser repetitiva y de distribución inestable, por lo que su aporte a la generalización de la alignment es limitado.}

\subsection{Dónde ocurre la alineación abierta}

La última diapositiva de contenido enumera participantes como AI2, HuggingFaceH4, Berkeley, LMSYS y NousResearch. La lista no pretende ser una clasificación completa, sino mostrar que el progreso abierto está distribuido en varios nodos: modelos, datos, trainer, evaluation y gestión de la comunidad; un ecosistema sano requiere que estos artifacts puedan combinarse entre sí y auditarse.

\lecturefigure{slide-076.jpg}{Principales comunidades, instituciones y tipos de artifact de la alignment abierta}{official deck, p.~76}

\begin{knowledgebox}{Lectura de la figura: la división del trabajo detrás de la lista de instituciones}
AI2 se centra en Tulu/OLMo y en la publicación de datos; HuggingFaceH4 reproduce recipes con rapidez; LMSYS aporta Arena y datos de conversaciones reales; Berkeley y los equipos de la comunidad exploran PPO y la evaluación; otros equipos generan instructions o preference data. La ventaja del ecosistema abierto es la exploración en paralelo; su debilidad es que la provenance, las licencias y los criterios de evaluación no están unificados.
\end{knowledgebox}

\teachervoice{La recomendación de investigación de 01:04:49--01:05:30 es: nadie puede seguir todos los avances, así que conviene seguir desarrollando habilidades y construir lo que uno considere importante. Para el ecosistema abierto, esto significa que aportar datos, evaluaciones o herramientas de alta calidad es igualmente investigación central, sin necesidad de perseguir solo la escala de los modelos.}

\subsection{Synthetic data y su relación con los datos humanos}

Los datos sintéticos permiten sortear parte de los problemas de derechos de autor, privacidad y costo de anotación, y también generar tokens a mayor escala; sin embargo, los sesgos del teacher model, los patrones repetitivos y la falta de intent de usuarios reales limitan la generalización. Una dirección más razonable es un pipeline mixto de seed humana, ampliación con modelos, filtrado automático y auditoría manual.

Sea $p_h$ la distribución de los datos reales y $p_s$ la distribución sintética; la distribución de entrenamiento mixta puede escribirse como
\[
p_{\lambda}=\lambda p_h+(1-\lambda)p_s,
\]
donde $\lambda$ controla la proporción de datos humanos. El $\lambda$ óptimo no es una constante fija: cuando cambian la calidad de los datos, el dominio de la tarea o la teacher strength, debe volver a elegirse mediante una held-out human evaluation.

\teachervoice{En la Q\&A, de 01:14:38--01:15:29, Nathan considera que la synthetic data es una de las pocas vías que permiten seguir generando más tokens, y cree que los modelos seguirán participando en el entrenamiento de otros modelos; pero aclara que esta dirección todavía está en una etapa muy temprana.}

\subsection{Resumen de la sección}

El protagonista del ecosistema moderno pasó de ser un modelo único a una artifact network. Llama 3 muestra la combinación de una base a gran escala con post-training en múltiples etapas; la brecha entre modelos abiertos y cerrados varía según la tarea; la preference data de alta calidad y la evaluación confiable siguen siendo recursos escasos. El progreso futuro podría venir de algoritmos unificados más simples, o bien de una mejor infraestructura de datos y de retroalimentación.
\section{Síntesis y ampliación}

La aportación más importante de esta clase no es clasificar DPO, PPO o QLoRA, sino ofrecer una historia interpretable de los sistemas abiertos de alineación: cada salto de los modelos dependió a la vez del base model, de la distribución de los datos, del acceso al entrenamiento, de la retroalimentación de la evaluación y de los artifacts publicados. Solo al entender esta cadena se puede juzgar qué cuello de botella resuelve realmente un método nuevo.

\subsection{Catorce conclusiones centrales}

Esta sección condensa el curso en conclusiones reutilizables; cada una conserva sus condiciones y su nivel de evidencia.

\begin{enumerate}
  \item RLHF es importante para el ChatGPT-like behavior, pero no es condición suficiente para un producto completo.
  \item El base model es la base de las capacidades; la alignment principalmente reordena y moldea la distribución de las interacciones.
  \item IFT, SFT, RLHF y DPO deben distinguirse por sus datos y su función de pérdida.
  \item Alpaca demostró que la synthetic instruction data permite adaptar la interfaz con rapidez.
  \item Vicuna/ShareGPT demostró que la prompt distribution de usuarios reales es sumamente valiosa.
  \item OpenAssistant muestra que los datos humanos abiertos tienen un valor de reutilización a largo plazo.
  \item QLoRA amplió el número de participantes, pero no resolvió automáticamente la estabilidad de RLHF.
  \item El debate entre safety y uncensored es, en esencia, una elección sobre el límite de las negativas y el control del usuario.
  \item Arena, AlpacaEval, MT-Bench y Leaderboard sirven a plazos de retroalimentación distintos.
  \item LLM-as-a-judge es rápido, pero tiene sesgos de length, position, self-preference y contamination.
  \item El objetivo de RLHF combina reward y KL constraint; PPO es solo una de las vías de implementación.
  \item DPO es sencillo y eficaz, pero difiere de PPO en la dinámica de entrenamiento y en el aprovechamiento de los datos.
  \item Zephyr/Tulu demostraron que DPO puede llevarse a la práctica; SteerLM/Starling demostraron que RL sigue siendo competitivo.
  \item En 2024, el cuello de botella persistente de la alignment abierta es la preference data diversa, legal y publicable.
\end{enumerate}

\subsection{Una tabla que articula el sistema abierto de alineación}

Para no quedarse solo con nombres de modelos, a continuación se resume, por componente del sistema, el problema que resolvió cada etapa y el riesgo residual.

\begin{table}[H]
\centering
\small
\begin{tabularx}{\textwidth}{L{0.16\textwidth} L{0.22\textwidth} X}
\toprule
Componente & Artifact representativo & Problemas aún sin resolver \\
\midrule
Base model & LLaMA/Llama 2/Llama 3 & Capacidades, licencia, provenance de los datos y costo de entrenamiento \\
IFT data & Alpaca, ShareGPT, UltraChat & Intent real, diversidad, privacidad y sesgo de plantilla \\
Human data & OpenAssistant/Arena logs & Costo, consistencia, consent y cobertura \\
Efficient tuning & LoRA/QLoRA & Conservación completa de las capacidades, RL stability y fusión para inference \\
Evaluation & Arena/AlpacaEval/MT-Bench/Leaderboard & Judge bias, contaminación, latencia y gaming \\
Preference optimization & PPO/DPO/SteerLM & Eficiencia de datos, estabilidad, reward hacking e interpretabilidad \\
Ecosystem & AI2, HF, LMSYS, community & Lineage, reproducibilidad, mantenimiento a largo plazo y gobernanza \\
\bottomrule
\end{tabularx}
\caption{Componentes de la alineación abierta, artifacts representativos y riesgos residuales.}
\end{table}

\subsection{Release checklist para la práctica}

La historia del curso puede convertirse en una checklist previa a la publicación de un modelo. El objetivo es que «usamos DPO» deje de sustituir a la evidencia completa.

\begin{lstlisting}[caption={Open aligned model release checklist}]
record_base_model_and_license()
publish_chat_template_and_training_objective()
document_instruction_and_preference_data_provenance()
report_compute_memory_and_random_seeds()
run_capability_safety_and_human_preference_evals()
disclose_judge_models_baselines_and_known_biases()
release_failure_examples_and_model_card()
\end{lstlisting}

\begin{importantbox}{Cinco filtros para juzgar un nuevo método de alignment}
\begin{enumerate}
  \item ¿Supera a un baseline fuerte con base/data/eval compartidos?
  \item ¿Produce un modelo competitivo descargable y reproducible, y no solo curvas de entrenamiento?
  \item ¿Conserva las base capabilities y reduce la failure rate real?
  \item ¿Se sostiene con varios judges y con preferencias humanas, y no en un único leaderboard?
  \item ¿La ganancia justifica los datos, el rollout, la memoria de GPU y la complejidad de ingeniería adicionales?
\end{enumerate}
\end{importantbox}

\subsection{Preguntas de autoevaluación}

Las siguientes preguntas sirven para comprobar si se domina la causalidad del sistema, y no solo la cronología de los releases.

\begin{enumerate}
  \item ¿Por qué RLHF es importante para ChatGPT y, aun así, no es condición suficiente?
  \item ¿Cómo se distinguen IFT y SFT en objetivo, datos e interfaz?
  \item ¿Por qué ShareGPT puede ser más valioso que una mayor cantidad de synthetic prompts?
  \item ¿Cuál es la diferencia esencial entre los weight differences y una LoRA update?
  \item ¿Qué memoria ahorra QLoRA y qué costos de RLHF no ahorra?
  \item ¿Para qué decisiones sirve cada uno de Arena, AlpacaEval, MT-Bench y Leaderboard?
  \item ¿Qué evita la KL penalty en el RLHF objective?
  \item ¿Qué representa el log-ratio policy/reference en la DPO loss?
  \item ¿Por qué Zephyr/Tulu no bastan por sí solos para demostrar que DPO siempre supera a PPO?
  \item ¿Por qué la preference data sintética puede ampliar la escala y, a la vez, estrechar la distribución?
\end{enumerate}

\subsection{Lecturas adicionales}

\begin{itemize}
  \item Ouyang et al., \emph{Training language models to follow instructions with human feedback}: InstructGPT y el pipeline de RLHF.
  \item Taori et al., \emph{Stanford Alpaca}; Wang et al., \emph{Self-Instruct}: la synthetic instruction data temprana.
  \item Dettmers et al., \emph{QLoRA}: base model cuantizado y fine-tuning de bajo rango.
  \item Rafailov et al., \emph{Direct Preference Optimization}: derivación y experimentos de DPO.
  \item Bai et al., \emph{Constitutional AI}; Köpf et al., \emph{OpenAssistant Conversations}: preferencias y datos humanos abiertos.
  \item Zheng et al., \emph{Chatbot Arena}: pairwise evaluation humana y el ecosistema de modelos abiertos.
\end{itemize}

\end{document}
